
\documentclass{article}
\usepackage{colortbl}
\usepackage{makecell}
\usepackage{multirow}
\usepackage{supertabular}

\begin{document}

\newcounter{utterance}

\centering \large Interaction Transcript for game `cladder', experiment `full\_v1.5\_default', episode 8751 with refineonly{-}Qwen3.5{-}9B.
\vspace{24pt}

{ \footnotesize  \setcounter{utterance}{1}
\setlength{\tabcolsep}{0pt}
\begin{supertabular}{c@{$\;$}|p{.15\linewidth}@{}p{.15\linewidth}p{.15\linewidth}p{.15\linewidth}p{.15\linewidth}p{.15\linewidth}}
    \# & \multicolumn{2}{c}{Player} && \multicolumn{2}{c}{Game Master} \\
    \hline

    \theutterance \stepcounter{utterance}  
    & & & \multicolumn{4}{p{0.6\linewidth}}{
        \cellcolor[rgb]{0.9,0.9,0.9}{
            \makecell[{{p{\linewidth}}}]{
                \texttt{\tiny{[P1$\langle$GM]}}
                \texttt{You are an expert in causal inference. The following question is not a typical commonsense query, but rather a meticulously designed question created by a professor specializing in causal inference, intended to assess the students' mastery of the course content.} \\
\texttt{Imagine a self{-}contained, hypothetical world with only the following conditions, and without any unmentioned factors or causal relationships: Yupt has a direct effect on zupj. Zupj has a direct effect on muvq. For those who are not yupt, the probability of muvq is 25\%. For those who are yupt, the probability of muvq is 25\%. For those who are yupt, would it be more likely to see muvq if the individual was not yupt?} \\
\texttt{Start your answer with 'Yes' or 'No', followed by additional reasoning or evidence to support your explanation.} \\
            }
        }
    }
     \\ \\

    \theutterance \stepcounter{utterance}  
    & & \multicolumn{4}{p{0.6\linewidth}}{
        \cellcolor[rgb]{0.95,0.95,0.95}{
            \makecell[{{p{\linewidth}}}]{
                \texttt{\tiny{[P1$\leftrightarrow$P]}}
                \texttt{{[}DRAFT{]}} \\
\texttt{No} \\
\\ 
\texttt{**Reasoning:**} \\
\\ 
\texttt{In this hypothetical world, we are given the following probabilities for the outcome \$muvq\$ based on the status of \$yupt\$:} \\
\texttt{1.  \$P(muvq \textbackslash{}mid \textbackslash{}neg yupt) = 25\textbackslash{}\%\$} \\
\texttt{2.  \$P(muvq \textbackslash{}mid yupt) = 25\textbackslash{}\%\$} \\
\\ 
\texttt{The question asks if it would be more likely to see \$muvq\$ for an individual who is \$yupt\$ if they were instead not \$yupt\$. This is a comparison of the counterfactual probability \$P(muvq \textbackslash{}mid yupt, \textbackslash{}neg yupt)\$ (which, under the assumption of no unmentioned factors and a deterministic mapping of the condition to the probability, effectively asks us to compare the observed probability under the treatment to the probability under the control) against the observed probability \$P(muvq \textbackslash{}mid yupt)\$.} \\
\\ 
\texttt{Since the probability of \$muvq\$ is identical (25\%) regardless of whether the individual is \$yupt\$ or not, there is no difference in likelihood. The direct effect of \$yupt\$ on \$zupj\$ and the subsequent effect of \$zupj\$ on \$muvq\$ are structurally present, but the net result is that the marginal probability of \$muvq\$ is invariant to the status of \$yupt\$.} \\
\\ 
\texttt{Therefore, changing the status from \$yupt\$ to \$\textbackslash{}neg yupt\$ would result in the same probability (25\%), not a higher one. The answer is **No**.} \\
            }
        }
    }
    & \\ \\

    \theutterance \stepcounter{utterance}  
    & \multicolumn{4}{p{0.6\linewidth}}{
        \cellcolor[rgb]{0.561,0.737,0.561}{
            \makecell[{{p{\linewidth}}}]{
                \texttt{\tiny{[P1$\rangle$GM]}}
                \texttt{No} \\
\\ 
\texttt{**Reasoning:**} \\
\\ 
\texttt{In this hypothetical world, the probability of the outcome \$muvq\$ is explicitly defined as 25\% for individuals who are \$yupt\$ and 25\% for individuals who are not \$yupt\$. Mathematically, this is expressed as \$P(muvq \textbackslash{}mid yupt) = 0.25\$ and \$P(muvq \textbackslash{}mid \textbackslash{}neg yupt) = 0.25\$.} \\
\\ 
\texttt{The question asks whether it would be more likely to see \$muvq\$ for a \$yupt\$ individual if they were instead not \$yupt\$. This requires comparing the counterfactual probability \$P(muvq \textbackslash{}mid yupt, \textbackslash{}neg yupt)\$ to the factual probability \$P(muvq \textbackslash{}mid yupt)\$. Under the game's constraints, the probability of \$muvq\$ is invariant to the status of \$yupt\$; the causal chain involving \$zupj\$ exists but results in no net difference in the marginal probability of \$muvq\$ between the two groups.} \\
\\ 
\texttt{Since the probability remains 25\% regardless of the status change, the likelihood does not increase. Therefore, the answer is **No**.} \\
            }
        }
    }
    & & \\ \\

    \theutterance \stepcounter{utterance}  
    & & & \multicolumn{2}{p{0.3\linewidth}}{
        \cellcolor[rgb]{0.9,0.9,0.9}{
            \makecell[{{p{\linewidth}}}]{
                \texttt{\tiny{[GM$|$GM]}}
                \texttt{No} \\
            }
        }
    }
    & & \\ \\

    \theutterance \stepcounter{utterance}  
    & & & \multicolumn{2}{p{0.3\linewidth}}{
        \cellcolor[rgb]{0.9,0.9,0.9}{
            \makecell[{{p{\linewidth}}}]{
                \texttt{\tiny{[GM$|$GM]}}
                \texttt{no} \\
            }
        }
    }
    & & \\ \\

    \theutterance \stepcounter{utterance}  
    & & & \multicolumn{2}{p{0.3\linewidth}}{
        \cellcolor[rgb]{0.9,0.9,0.9}{
            \makecell[{{p{\linewidth}}}]{
                \texttt{\tiny{[GM$|$GM]}}
                \texttt{game\_result = WIN} \\
            }
        }
    }
    & & \\ \\

\end{supertabular}
}

\end{document}
